%% ============================================================
%% Chapter 5 — Operational Intelligence
%% Source: operational-agent-intelligence.tex
%% Two-column body layout via multicols; full-width figures between blocks.
%% ============================================================

%% Objects re-used from Chapters 3–4 (providecommand avoids redefinition errors)
\providecommand{\receiver}{\mathsf{R}}
\providecommand{\agent}{\mathcal{A}}
\providecommand{\Cell}{\mathcal{C}}
\providecommand{\Sfloor}{S_{\flat}}
\providecommand{\Mode}{\mathfrak{M}}
\providecommand{\Mfld}{\mathcal{M}}
\providecommand{\Pot}{\Phi}
%% New in ch05
\newcommand{\KFW}{\mathcal{K}}
\newcommand{\Iindex}{\mathcal{I}}
\newcommand{\Acon}{A_{\mathrm{con}}}
\newcommand{\Aact}{A_{\mathrm{act}}}
\newcommand{\Tcon}{T_{\mathrm{con}}}
\newcommand{\hbarA}{\hbar_{\!\mathcal{A}}}

\chapter{Operational Intelligence: Admissibility, Cycle Necessity, and the Intelligence Index}
\label{ch:intelligence}

\papernote{\citep{sachikonye2024_intelligence}}

\begin{chapterabstract}
A substrate-neutral operational definition of intelligence is developed for any bounded
agent. The definition rests on four conditions: cell-disjoint extension of the agent's
knowledge framework, novel-cell construction, floor reduction on novel cells, and
synchronisation feasibility with other agents. From the receiver calculus established in
Chapter~\ref{ch:purpose} --- the residual semantic distance $S(\receiver, x; \Cell)$
and its irreducible floor $\beta$ --- seven theorems fix what intelligence is, when it
can occur, how to measure it, why it is bounded, and how its failures classify.
Extension is admissible iff cell-disjoint~(T1); extension occurs only during phases in
which the agent's action-space dispersion is suspended, by a Heisenberg-type uncertainty
inequality on receiver dispersions~(T2); intelligence requires alternation between
construction and action phases~(T3); the intelligence index $\Iindex(\agent)$, computed
from per-phase activity ratios, phase-time fractions, and perceptual coupling, is
positive iff the agent is intelligent~(T4); collective intelligence is precisely
synchronisation of agents on novel cells via aperture-sharing isomorphisms~(T5); every
realised agent has finite intelligence index --- there is no maximally intelligent
agent~(T6); cycle failures admit exactly six distinct phenotypes, each with a closed-form
signature in $\Iindex$ and its component factors~(T7). Thirty numerical experiments
verify all claims to machine precision.
\end{chapterabstract}

%% ============================================================
%% Block 1 — The substrate-neutral question, axioms, agent
%% ============================================================
\begin{multicols}{2}

\section{The substrate-neutral question}
\label{sec:ch05-intro}

The question \emph{what does it mean to be an intelligent agent?} has been posed for as
long as agents have been studied. Most answers import substrate-specific assumptions ---
biological neural circuits, symbolic computation, observable behaviour, evolutionary
fitness --- that restrict the answer to one kind of agent. The consequence is one notion
of intelligence for humans, another for animals, another for machines, and no unified
framework across substrates.

This chapter develops a substrate-neutral definition: a single mathematical statement
that applies uniformly to biological, computational, and organisational agents, and that
yields the same operational test regardless of substrate. The definition is built
entirely from the receiver calculus established in Chapters~\ref{ch:purpose}
and~\ref{ch:cohesion}: the $S$-functional, the receiver floor~$\beta$, the action-cell
tolerance~$\tau$, and the Lagrangian agent on the state manifold~$\Mfld$.

\subsection{What this definition rules out}

The definition is exclusionary by design. Several common notions of intelligence are
formally ruled out.

\textbf{Intelligence as success rate.} Floor positivity (Theorem~\ref{thm:floor},
Chapter~\ref{ch:purpose}) guarantees every bounded agent fails some fraction of tasks.
A definition equating intelligence with high success rate is incompatible with floor
positivity, hence with bounded agency.

\textbf{Intelligence as common knowledge.} Common-Cell Convergence
(Theorem~\ref{thm:ccc}, Chapter~\ref{ch:purpose}) shows that agents with disjoint
representations can coordinate on shared cells without shared knowledge. Coordination is
not intelligence; intelligence does not require shared concepts.

\textbf{Intelligence as point-correctness.} Cell-truth (Theorem~\ref{thm:cell-truth},
Chapter~\ref{ch:purpose}) says output-equivalence holds at the cell level. Distinguishing
point outputs within a cell is forbidden by the receiver floor; doing so is a category
error, not an achievement.

\textbf{Intelligence as synchronisation.} A fully synchronised collective is not
individually intelligent. The phase-locked regime ($R \ge 0.95$) established in
Chapter~\ref{ch:cohesion} is a coordination property, not an intelligence certificate.

\textbf{Intelligence as meta-methodology.} Selecting which methodology to apply is
methodology-routing. An agent is intelligent only when the \emph{substrate of selection
itself} can grow. The remainder of this chapter develops the consequences of that
restriction.

\section{Axiomatic foundations and the bounded agent}
\label{sec:ch05-axioms}

The three axioms governing all agents in this framework were introduced in
Chapter~\ref{ch:cohesion}. They are recalled here for self-containment.

\begin{axiom}[Bounded Phase Space]\label{ax:ch05-bounded}
Every agent occupies a phase space $\Omega$ of finite measure: $\mu(\Omega) < \infty$.
\end{axiom}

\begin{axiom}[No Null State]\label{ax:ch05-null}
At every instant, an agent occupies exactly one categorical state. There is no null
state and no inter-state void.
\end{axiom}

\begin{axiom}[Finite Resolution]\label{ax:ch05-resolution}
Any observation, internal or external, has finite resolution $\delta > 0$. States
separated by less than $\delta$ are observationally indistinguishable.
\end{axiom}

\begin{remark}[Universality of the axioms]
All three axioms hold across substrates: biological agents face finite metabolic budgets,
perpetual oscillation, and finite sensor resolution; computational agents face finite
memory, perpetual instruction execution, and finite floating-point precision;
organisational agents face finite resources, perpetual operation, and finite
decision-making bandwidth. No substrate-specific assumption is made.
\end{remark}

\subsection{Outcome space, action-cells, and the receiver}

The outcome space $(\mathcal{X}, d, \Sigma)$, action map $A: \mathcal{X} \to \mathcal{Y}$,
action-cells $\Cell_y = A^{-1}(\{y\})$ with tolerance
$\tau(\Cell_y) = \sup_{x \in \mathcal{X}} \inf_{c \in \Cell_y} d(x,c)$, and receiver
$\receiver = (\KFW, \mathcal{D}, \Pi, \beta)$ with knowledge framework $\KFW$, decoder
$\mathcal{D}$, candidate-projection $\Pi$, and floor $\beta > 0$ were all established in
Chapter~\ref{ch:purpose}. The $S$-functional and methodology are recalled for
completeness.

\begin{definition}[$S$-functional, recalled]\label{def:ch05-S}
For receiver $\receiver$, state $x \in \mathcal{X}$, and action-cell $\Cell$,
\[
S(\receiver, x; \Cell) := \inf_{x' \in \Pi(\mathcal{D}(x))} d(x', \Cell) + \beta.
\]
The floor is $\Sfloor(\receiver) := \beta > 0$.
\end{definition}

\begin{definition}[Methodology]\label{def:ch05-method}
A \emph{methodology} on receiver $\receiver$ is a triple
$\mathsf{M} = (T, \kappa, \sigma)$ where $T: \mathcal{X} \times [0, \Sigma] \to
[0, \Sigma]$ is a per-iteration update operator, $\kappa \in [0,1)$ is the contraction
factor, and $\sigma \in [0, \Sigma]$ is the per-iteration dispersion. The intrinsic
floor is $\Sfloor(\mathsf{M}) = \sigma\kappa/(1-\kappa)$.
\end{definition}

\section{The agent: state manifold, trajectory, and aperture}
\label{sec:ch05-agent}

The Lagrangian agent structure developed in Chapter~\ref{ch:cohesion} is the operational
basis for the intelligence theorems that follow. Key objects are recalled.

\begin{definition}[State manifold, recalled]\label{def:ch05-manifold}
The agent's state manifold is the compact set
\[
\Mfld = [0,1] \times [0, 2\pi^2] \times [0,1]^3,
\]
parametrised by generalised coordinates $q = (R, \sigma^2, S_k, S_t, S_e)$, where
$R \in [0,1]$ is internal coherence, $\sigma^2$ is phase variance, and
$(S_k, S_t, S_e)$ are $S$-entropy coordinates.
\end{definition}

\begin{definition}[Agent state triple, recalled]\label{def:ch05-triple}
The complete state of an agent is the triple
$\Mode = (\gamma, \Gamma_f, M)$,
where $\gamma: [0,T] \to \Mfld$ is the trajectory, $\Gamma_f = \gamma(T)$ is the
terminus, and $M = \int_0^T \dot{H}^+\,\mathrm{d}t$ is the memory.
\end{definition}

\begin{definition}[Agent, recalled]\label{def:ch05-agent}
An \emph{agent} is the tuple
$\agent = (q, \gamma, \Gamma_f, M, \Aper, \Pot, \mathrm{regime})$,
where $q \in \Mfld$ is the current state, $\Pot: \Mfld \to \mathbb{R}$ is the partition
potential, and $\mathrm{regime} \in \{\textsc{turbulent}, \textsc{aperture},
\textsc{cascade}, \textsc{coherent}, \textsc{phase-locked}\}$ is classified by $R$ as
in Chapter~\ref{ch:cohesion}.
\end{definition}

Apertures were shown to perform zero thermodynamic work
(Theorem~\ref{thm:zero-work}, Chapter~\ref{ch:cohesion}): they constrain accessible
states without expending energy. This zero-work property is what makes
apertures a substrate-neutral information-processing primitive.

\end{multicols}

%% ============================================================
%% Panel 1 — Foundations of the Intelligent Agent
%% ============================================================
\begin{figure}[!htbp]
    \centering
    \includegraphics[width=\textwidth]{figures/operational-intelligence/panel_1.png}
    \caption{\textbf{Foundations of the Intelligent Agent.}
    (\textbf{A})~Floor positivity across receivers: the floor $\beta(\receiver)$ on the
    linear $y$-axis versus knowledge framework size $|\KFW|$ on the log $x$-axis ($2$ to
    $1024$). The floor decreases monotonically with $|\KFW|$ but remains strictly positive
    at every tested size; the crimson dashed reference at $\beta = 0$ is forbidden by the
    Floor Theorem (Chapter~\ref{ch:purpose}). (\textbf{B})~Cell-disjointness criterion:
    bar chart of cell-overlap counts for $50$ random aperture pairs. Seagreen bars (zero
    overlap) indicate cell-disjoint pairs admissible under Theorem~\ref{thm:ch05-t1};
    orange bars indicate inadmissible pairs that violate cell-uniqueness.
    (\textbf{C})~Receiver Uncertainty Principle: scatter of $\sigma_Y$ versus $\sigma_K$
    for $100$ random methodologies. All measured products lie at or above the bound
    $\sigma_K \sigma_Y = \hbarA$; zero violations (Theorem~\ref{thm:ch05-uncertainty}).
    (\textbf{D})~Three-dimensional surface of $\sigma_K \cdot \sigma_Y$ over
    $(\sigma_K, \sigma_Y) \in [0.2,5]^2$, viridis colourmap. The crimson contour at
    $\hbarA = \beta\tau$ traces the saturating curve; the surface lies above it
    everywhere, geometrically realising the agent phase-space cell of area $\hbarA$.}
    \label{fig:ch05-panel1-foundations}
\end{figure}

%% ============================================================
%% Block 2 — Dispersions and the intelligence cycle
%% ============================================================
\begin{multicols}{2}

\section{Conjugate dispersions and the Receiver Uncertainty Principle}
\label{sec:ch05-dispersions}

\subsection{Knowledge and action dispersions}

The methodology's per-iteration dispersion $\sigma$ (Definition~\ref{def:ch05-method})
decomposes along two conjugate axes corresponding to the agent's knowledge framework and
its action space. This decomposition is the starting point for the uncertainty theorem.

\begin{definition}[Knowledge dispersion]\label{def:ch05-dispK}
For agent $\agent$ with methodology $\mathsf{M}$ at state $x$, let
$\mathsf{M}^{(1)}(x), \mathsf{M}^{(2)}(x) \in \KFW$ be two independent applications
of one iteration of $\mathsf{M}$ at $x$. The \emph{knowledge dispersion} is
\[
\sigma_K(\mathsf{M}, x) :=
  \mathbb{E}\bigl[d_K\bigl(\mathsf{M}^{(1)}(x), \mathsf{M}^{(2)}(x)\bigr)\bigr],
\]
where $d_K$ is the metric on $\KFW$ induced by the decoder $\mathcal{D}$. The overall
knowledge dispersion is $\sigma_K(\mathsf{M}) = \sup_x \sigma_K(\mathsf{M}, x)$.
\end{definition}

\begin{definition}[Action dispersion]\label{def:ch05-dispY}
With $A$ the action map and $d_Y$ a metric on $\mathcal{Y}$,
\[
\sigma_Y(\mathsf{M}, x) :=
  \mathbb{E}\bigl[d_Y\bigl(A(\mathsf{M}^{(1)}(x)),
  A(\mathsf{M}^{(2)}(x))\bigr)\bigr],
\]
with $\sigma_Y(\mathsf{M}) = \sup_x \sigma_Y(\mathsf{M}, x)$.
\end{definition}

The two dispersions are conjugate: reducing uncertainty in the action space (committing
to an output) forces expansion in the knowledge space (opening representational
flexibility), and vice versa. This conjugacy is the formal content of not knowing what
you do not know.

\subsection{The agent phase-space cell}

\begin{definition}[Agent phase-space cell]\label{def:ch05-phasecell}
The \emph{phase-space cell} of agent $\agent$ on action-cell $\Cell$ of tolerance $\tau$
is the rectangle $[\beta] \times [\tau] \subset \KFW \times \mathcal{Y}$, with area
\[
\hbarA := \beta \cdot \tau.
\]
\end{definition}

The quantity $\hbarA$ is the agent-level analogue of Planck's constant. Just as $\hbar$
measures the irreducible quantum of physical action, $\hbarA$ measures the irreducible
quantum of the agent's joint knowledge--action resolution. The analogy with the
Heisenberg uncertainty principle~\citep{heisenberg1927,robertson1929} and the
Cram\'er--Rao bound~\citep{rao1945,cramer1946} is made precise in the following theorem.

\subsection{The Receiver Uncertainty Principle}

\begin{theorem}[Receiver Uncertainty Principle]\label{thm:ch05-uncertainty}
For any methodology $\mathsf{M}$ on agent $\agent$ with action-cell tolerance $\tau$,
\[
\sigma_K(\mathsf{M}) \cdot \sigma_Y(\mathsf{M}) \ge \hbarA = \beta \cdot \tau.
\]
\end{theorem}

\begin{proof}
Two independent samples of one iteration of $\mathsf{M}$ at fixed input are
agent-distinguishable only when they fall into different phase-space cells, i.e., when
$d_K > \beta$ \emph{and} $d_Y > \tau$. By Markov's inequality,
\[
\Pr[d_K \le \beta] \ge 1 - \sigma_K/\beta, \quad
\Pr[d_Y \le \tau] \ge 1 - \sigma_Y/\tau.
\]
For the methodology to deliver at least one bit per pair of samples, the same-cell
probability must be strictly less than $1$. Both factors must therefore be strictly less
than $1$, which is equivalent to $\sigma_K \sigma_Y \ge \beta\tau$ at the saturating
limit.
\end{proof}

\begin{remark}[Informational content]
High action certainty (small $\sigma_Y$, near-deterministic output) forces large
representational uncertainty (large $\sigma_K$, many options in $\KFW$). This is the
operational form of the observation that high output commitment comes at the cost of
representational flexibility: the agent cannot simultaneously be decided in what it
does and undecided in what it knows.
\end{remark}

\section{Knowledge framework extension and the intelligence cycle}
\label{sec:ch05-cycle}

\subsection{Extension and admissibility}

\begin{definition}[Knowledge framework extension]\label{def:ch05-extension}
A \emph{knowledge framework extension} is a relation
$\KFW_t \subseteq \KFW_{t+\Delta t}$ between the agent's knowledge framework at
successive times. The extension is \emph{strict} if
$\KFW_{t+\Delta t} \setminus \KFW_t \neq \varnothing$, and \emph{trivial} otherwise.
\end{definition}

\begin{theorem}[Extension Admissibility, T1]\label{thm:ch05-t1}
A strict extension $\KFW_{t+\Delta t} \supsetneq \KFW_t$ is admissible if and only if
every newly added categorical distinction
$k^{\mathrm{new}} \in \KFW_{t+\Delta t} \setminus \KFW_t$ is cell-disjoint from every
distinction in $\KFW_t$: for every $k^{\mathrm{old}} \in \KFW_t$, the corresponding
action-cells satisfy $\Cell^{\mathrm{new}} \cap \Cell^{\mathrm{old}}$ has zero tolerance.
\end{theorem}

\begin{proof}
Cell-truth (Chapter~\ref{ch:purpose}, Theorem~\ref{thm:cell-truth}) requires every state
to be assigned to exactly one cell. If $\Cell^{\mathrm{new}} \cap \Cell^{\mathrm{old}}$
has positive tolerance, states in the intersection are assigned to both cells, violating
cell-uniqueness and invalidating prior cell-attainment. This establishes necessity.

For sufficiency: if every newly added distinction is cell-disjoint from every prior
distinction, the partition refinement preserves the prior partition. Every prior cell
$\Cell^{\mathrm{old}}$ remains an action-cell (it is the union of its refined sub-cells),
and prior cell-attainment is preserved by inclusion.
\end{proof}

\begin{remark}[Accumulation without collision]
Admissible extension is the formal content of learning without forgetting: new
representations must be orthogonal (cell-disjoint) to old ones, not merely different
parameter values within the same representational bucket. This rules out pseudolearning
--- re-encoding existing distinctions as new symbols --- as a mechanism of intelligence.
\end{remark}

\subsection{Construction and action phases}

\begin{definition}[Construction phase, action phase]\label{def:ch05-phases}
A \emph{construction phase} of agent $\agent$ is a temporal interval during which
$\sigma_Y \to 0$ (the agent commits to no actions). An \emph{action phase} is a temporal
interval during which $\sigma_K \to 0$ (the agent commits to a single decoded
representation).
\end{definition}

The construction--action distinction maps onto the Lagrangian regime structure of
Chapter~\ref{ch:cohesion}: a construction phase corresponds to traversal of
high-$\sigma_K$ regions of $\Mfld$ (turbulent and aperture regimes), while an action
phase corresponds to traversal of high-$\sigma_Y$ regions (cascade through phase-locked).
The five-regime structure is the dynamical substrate for phase alternation.

\begin{theorem}[Construction-Phase Necessity, T2]\label{thm:ch05-t2}
Knowledge framework extension can occur only during construction phases.
\end{theorem}

\begin{proof}
To add new categorical distinctions to $\KFW$, the per-iteration update must explore
representations not present in $\KFW_t$. Fundamentally new representations require
$\sigma_K > \beta$ (distinctions within $\beta$ of existing ones are not new by
cell-truth). At the extreme, $\sigma_K \to \infty$ in the saturating limit.

By the Receiver Uncertainty Principle (Theorem~\ref{thm:ch05-uncertainty}),
$\sigma_K \cdot \sigma_Y \ge \hbarA$. For $\sigma_K \to \infty$, the inequality forces
$\sigma_Y \to 0$ --- the definition of a construction phase.

Conversely, any phase with $\sigma_Y > 0$ bounds $\sigma_K \le \hbarA/\sigma_Y < \infty$,
restricting the agent to representations within finite dispersion of $\KFW_t$. These are
parameter updates within the existing framework, not framework extension.
\end{proof}

\subsection{Cycle necessity}

\begin{theorem}[Cycle Necessity, T3]\label{thm:ch05-t3}
An agent operating exclusively in either construction phase or action phase cannot be
intelligent. Intelligence requires alternation.
\end{theorem}

\begin{proof}
The intelligence definition (Definition~\ref{def:ch05-intelligence}) requires both novel
cell construction (condition~(ii)) and floor reduction on novel cells (condition~(iii)).

\emph{Exclusive construction phase}: $\sigma_Y = 0$ throughout. The agent commits to no
actions. Floor reduction requires the candidate-projection to attain the target cell, and
attainment is verified through action commitment. Without commitment, condition~(iii) is
unverifiable.

\emph{Exclusive action phase}: $\sigma_Y > 0$ throughout. By
Theorem~\ref{thm:ch05-t2}, no extension occurs. The knowledge framework is fixed, so
condition~(ii) is violated at every step.

Both phases are required, and by Theorem~\ref{thm:ch05-t2} they cannot coexist.
Therefore intelligence requires alternation.
\end{proof}

\end{multicols}

%% ============================================================
%% Panel 2 — The Intelligence Cycle
%% ============================================================
\begin{figure}[!htbp]
    \centering
    \includegraphics[width=\textwidth]{figures/operational-intelligence/panel_2.png}
    \caption{\textbf{The Intelligence Cycle: Construction and Action Phases.}
    (\textbf{A})~Construction-phase necessity: required $\sigma_K$ versus $\sigma_Y$ on
    log--log axes for $50$ values. As $\sigma_Y \to 0$ the required $\sigma_K$ diverges,
    confirming Theorem~\ref{thm:ch05-t2}: knowledge framework extension can occur only
    during phases in which action commitment is suspended. (\textbf{B})~Phase alternation
    timeline: $\sigma_K$ (steelblue) and $\sigma_Y$ (orange) versus time ($0$ to $5$
    cycles). The signals are anti-phase, confirming Theorem~\ref{thm:ch05-t3}: the two
    phases are mutually exclusive. (\textbf{C})~Healthy activity ratio: intelligence proxy
    versus $\Acon/\Aact$; the peak at $\sqrt{0.95} \approx 0.975$ (crimson dashed)
    corresponds to the healthy reference operating ratio. (\textbf{D})~Cycle trajectory in
    $(\sigma_K, \sigma_Y, t)$ space: the closed orbit never enters the
    simultaneously-high quadrant, geometrically realising the mutual-exclusion property
    of Theorems~\ref{thm:ch05-t2} and~\ref{thm:ch05-t3}.}
    \label{fig:ch05-panel2-cycle}
\end{figure}

%% ============================================================
%% Block 3 — Intelligence index and floor
%% ============================================================
\begin{multicols}{2}

\section{Operational measurement: the intelligence index}
\label{sec:ch05-index}

\subsection{The four-condition definition}

\begin{definition}[Intelligence]\label{def:ch05-intelligence}
An agent $\agent$ is \emph{intelligent} iff there exist times $t_1 < t_2 < \cdots$ and
corresponding strict knowledge framework extensions
$\KFW_{t_i} \supsetneq \KFW_{t_{i-1}}$ such that, at each step:
\begin{enumerate}[label=(\roman*)]
\item \emph{Cell-disjoint extension}: Theorem~\ref{thm:ch05-t1} is satisfied at $t_i$;
\item \emph{Novel cell construction}: at $t_i$ there exists an action-cell $\Cell_i$
  that did not exist as a positive-tolerance cell under the prior framework
  $\KFW_{t_{i-1}}$;
\item \emph{Floor reduction}: $\Sfloor(\agent, \Cell_i; t_i) < \tau(\Cell_i)$;
\item \emph{Synchronisation feasibility}: there exists a class of agents with which
  $\agent$ could phase-lock on $\Cell_i$.
\end{enumerate}
\end{definition}

Condition~(i) forbids pseudolearning. Condition~(ii) requires genuinely new semantic
content. Condition~(iii) requires actionability: the agent must be able to attain the
novel cell, not merely represent it. Condition~(iv) opens the agent's intelligence to
social verification: a novel cell that no other agent could ever recognise is not
intelligence but idiosyncrasy.

\subsection{The intelligence index}

\begin{theorem}[Intelligence Index, T4]\label{thm:ch05-t4}
Define the \emph{intelligence cycle index}
\[
\Iindex(\agent) :=
  \frac{\Acon(\agent)}{\Aact(\agent)} \cdot
  \frac{\Tcon(\agent)}{\Tcon^{\mathrm{ref}}} \cdot
  \kappa_\agent,
\]
where $\Acon, \Aact$ are activity rates during construction and action phases
respectively, $\Tcon$ is the construction-phase duration per unit total time,
$\Tcon^{\mathrm{ref}}$ is a substrate-specific reference duration, and $\kappa_\agent$
is the perceptual coupling coefficient. Then $\Iindex(\agent) > 0$ is necessary for
intelligence; $\Iindex(\agent) \approx 1$ is the healthy operating point.
\end{theorem}

\begin{proof}
By Theorem~\ref{thm:ch05-t3}, intelligence requires both phases, so $\Tcon > 0$ and
$T_{\mathrm{act}} > 0$. By Theorem~\ref{thm:ch05-t2}, construction-phase activity must
reach the agent's full activity budget for $\sigma_K$ to be large enough to construct
novel cells satisfying condition~(ii).

The factor $\Acon/\Aact$ measures whether construction-phase activity is at full budget.
The factor $\Tcon/\Tcon^{\mathrm{ref}}$ measures whether the agent is allocated
sufficient construction time. The factor $\kappa_\agent$ measures whether the agent has
perceptual access to novel cells, required for testing condition~(iv).

If any factor is zero, at least one condition fails:
\begin{itemize}
\item $\Acon = 0$: no construction-phase activity, so no extension, so~(ii) fails.
\item $\Tcon = 0$: no construction time, so~(ii) fails.
\item $\kappa_\agent = 0$: no perceptual coupling, so~(iv) cannot be tested.
\end{itemize}
Conversely, all factors positive means the agent satisfies the necessary metabolic and
temporal preconditions for the intelligence cycle. The product is therefore positive iff
the agent can in principle satisfy the intelligence definition.

The reference value $\Iindex \approx 1$ corresponds to the agent operating at its full
activity budget with intact construction-phase duration and perceptual coupling. The
activity ratio target $\Acon/\Aact \approx \sqrt{0.95}$ follows from the critical
Kuramoto threshold $R \ge 0.95$ for phase-locking
(Chapter~\ref{ch:cohesion}, Theorem~\ref{thm:sync}): the construction-to-action ratio
that places the agent at the critical coupling boundary.
\end{proof}

\section{The intelligence floor}
\label{sec:ch05-floor}

\begin{theorem}[Intelligence Floor, T6]\label{thm:ch05-t6}
No realised agent achieves $\Iindex(\agent) = \infty$. Every realised intelligence is
bounded above.
\end{theorem}

\begin{proof}
By the Floor Theorem (Chapter~\ref{ch:purpose}), every bounded receiver has $\beta > 0$.
By the Receiver Uncertainty Principle, $\hbarA = \beta\tau > 0$ for any cell of positive
tolerance. By Theorem~\ref{thm:ch05-t2}, the construction phase requires $\sigma_Y \to 0$
and $\sigma_K$ bounded above by $\hbarA/\sigma_Y$ for any residual $\sigma_Y > 0$ in
the transition.

The activity rate $\Acon$ is bounded above by the agent's total activity budget. The
duration $\Tcon$ is bounded above by total available time. The perceptual coupling
$\kappa_\agent$ is bounded above by physical or architectural constraints.

Each factor in $\Iindex$ is bounded above; their product is bounded. The supremum
$\Iindex \to \infty$ would require $\beta \to 0$, forbidden by floor positivity.
Therefore no maximally intelligent agent can be realised.
\end{proof}

\begin{corollary}[Intelligence is degree]\label{cor:ch05-degree}
Intelligence is a degree on a bounded continuum. The relevant question is not
``intelligent vs.\ unintelligent'' but ``where on the continuum.''
\end{corollary}

\begin{remark}[The G\"{o}delian residue]
The floor $\beta > 0$ is the operational analogue of G\"{o}del's incompleteness
results~\citep{godel1931,tarski1933}: a strictly positive lower bound on representational
granularity from which no finite system can escape. The G\"{o}delian residue theorem
(Chapter~\ref{ch:purpose}, Theorem~\ref{thm:resid-floor}) identifies this residue as
exactly $\beta$. The intelligence cycle is the dynamical mechanism by which the residue
is reduced on novel cells via construction--action alternation --- but never to zero,
since $\beta > 0$ at all finite times.
\end{remark}

\end{multicols}

%% ============================================================
%% Panel 3 — The Intelligence Index
%% ============================================================
\begin{figure}[!htbp]
    \centering
    \includegraphics[width=\textwidth]{figures/operational-intelligence/panel_3.png}
    \caption{\textbf{The Intelligence Index.}
    (\textbf{A})~Index factor decomposition for four agent types: healthy (seagreen),
    construction-deficient P1 (crimson), construction-deprived P3 (orange), and
    perceptually decoupled P4 (purple). Each phenotype shows depression in exactly one
    factor with the others preserved --- the distinctive signature predicted by
    Theorem~\ref{thm:ch05-t7}. (\textbf{B})~Index sensitivity to perceptual coupling
    $\kappa$: the relationship is exactly linear, as predicted by
    Theorem~\ref{thm:ch05-t4} ($\Iindex$ is a product of independent factors). The
    crimson dashed reference at $\sqrt{0.95}$ is the healthy operating level.
    (\textbf{C})~Population distribution of normalised $\Iindex$ across $100$ simulated
    agents: centred tightly around $1$, demonstrating stability of the healthy operating
    point. (\textbf{D})~Three-dimensional surface of $\Iindex$ over
    $(\Acon/\Aact, \Tcon/\Tcon^{\mathrm{ref}})$ at $\kappa = 1$: bilinear, monotone in
    each factor, and bounded above by the agent's metabolic and temporal constraints
    (Theorem~\ref{thm:ch05-t6}).}
    \label{fig:ch05-panel3-index}
\end{figure}

%% ============================================================
%% Block 4 — Collective intelligence
%% ============================================================
\begin{multicols}{2}

\section{Collective intelligence}
\label{sec:ch05-collective}

\subsection{Federations and aperture-sharing}

\begin{definition}[Federation]\label{def:ch05-federation}
A \emph{federation} of agents is a tuple
$E = (\agent_1, \ldots, \agent_n)$ together with inter-agent coupling $K_{ij} \ge 0$
and natural-frequency vector $\omega = (\omega_1, \ldots, \omega_n)$.
\end{definition}

Federations are the multi-agent ensembles of Chapters~\ref{ch:purpose}
and~\ref{ch:cohesion}, now equipped with the additional structure required for collective
intelligence. The coupling--frequency structure is the Kuramoto interaction model whose
synchronisation threshold was proved in Chapter~\ref{ch:cohesion} (Theorem~\ref{thm:sync}).

\begin{definition}[Aperture-sharing isomorphism]\label{def:ch05-apsiso}
Two agents $\agent_i, \agent_j$ \emph{share apertures} on cell $\Cell$ if there exists
a structure-preserving isomorphism $\phi_{ij}: \KFW_i \to \KFW_j$ such that
$\mathcal{D}_j(x) = \phi_{ij}(\mathcal{D}_i(x))$ for all $x$ in the relevant domain,
and the corresponding candidate-projections identify under $\phi_{ij}$.
\end{definition}

Aperture-sharing is strictly stronger than common-cell attainment
(Chapter~\ref{ch:purpose}, Theorem~\ref{thm:ccc}). Common-Cell Convergence says
representation-disjoint agents can attain the same cell without sharing representations.
Aperture-sharing requires a decoder isomorphism: the agents' epistemic structures must
be mutually recognisable, not merely behaviourally aligned.

\subsection{Collective intelligence theorem}

\begin{theorem}[Collective Intelligence, T5]\label{thm:ch05-t5}
A federation $E$ is collectively intelligent iff for each agent's novel cell $\Cell_i$,
there exists an aperture-sharing isomorphism between every pair of agents mapping
$\Cell_i$ to its corresponding $\Cell_j$.
\end{theorem}

\begin{proof}
\emph{Sufficiency}: under the isomorphism condition, agent $\agent_j$ phase-locks with
$\agent_i$ on $\Cell_i$ via $\phi_{ij}$. The inter-agent partition lag vanishes by the
Synchronisation Theorem (Chapter~\ref{ch:cohesion}, Theorem~\ref{thm:sync}): when
agents become decoder-isomorphic, partition lag is zero and coordination friction
vanishes. The federation acts as a single collective agent with knowledge framework
$\bigcup_i \KFW_i$, satisfying conditions~(i)--(iv) of
Definition~\ref{def:ch05-intelligence} by construction.

\emph{Necessity}: condition~(iv) of Definition~\ref{def:ch05-intelligence} requires
synchronisation feasibility on novel cells. Without an aperture-sharing isomorphism,
agents' decoder--projection chains process the same input into incompatible categorical
states, and phase-locking is forbidden by Belief Incompatibility
(Chapter~\ref{ch:purpose}, Theorem~\ref{thm:belief-incompat}).
\end{proof}

\begin{remark}[Collective intelligence as partition extinction]
\label{rem:ch05-coll-pe}
The collective intelligence condition is the agent-level analogue of partition extinction
in physical systems (Chapter~\ref{ch:cohesion}). Just as paired electrons becoming
indistinguishable eliminates electrical resistance, decoder-isomorphic agents becoming
categorically indistinguishable eliminates coordination friction. The transition is
discontinuous: as with the Cooper-pair analogy (Remark~\ref{rem:cooper},
Chapter~\ref{ch:cohesion}), intra-collective phase-locking is binary once
$K \ge K_c^{\mathrm{ens}}$.
\end{remark}

\end{multicols}

%% ============================================================
%% Panel 4 — Collective Intelligence and Federations
%% ============================================================
\begin{figure}[!htbp]
    \centering
    \includegraphics[width=\textwidth]{figures/operational-intelligence/panel_4.png}
    \caption{\textbf{Collective Intelligence and Federations.}
    (\textbf{A})~Federation index growth: composite $\Iindex$ versus federation size $n$
    ($1$ to $20$) for four individual levels $\{0.3, 0.5, 0.7, 0.9\}$. Each curve
    approaches $1$ as $n$ grows, at a rate controlled by the individual level, confirming
    Theorem~\ref{thm:ch05-t5}. (\textbf{B})~Aperture-sharing matrix for $8$ agents ($4$
    shared, $4$ disjoint): the block structure confirms that only structurally compatible
    apertures support phase-locking. (\textbf{C})~Phase-locking gain: composite $\Iindex$
    (steelblue) and theoretical upper bound $1-(1-\Iindex_{\mathrm{indiv}})^n$ (crimson
    dashed) coincide, confirming the parallel-composition upper bound.
    (\textbf{D})~Composite-index surface over $(n_{\mathrm{agents}},
    \Iindex_{\mathrm{indiv}})$: saturates near $1$ for large $n$ and high individual
    levels; at small $n$ and low individual levels the composite remains low, showing
    diminishing returns of adding agents to an already high-intelligence federation.}
    \label{fig:ch05-panel4-collective}
\end{figure}

%% ============================================================
%% Block 5 — Failure-mode classification
%% ============================================================
\begin{multicols}{2}

\section{Failure-mode classification}
\label{sec:ch05-phenotypes}

\begin{theorem}[Failure-Mode Classification, T7]\label{thm:ch05-t7}
Cycle failures of an intelligent agent admit exactly six distinct phenotypes,
distinguishable through $\Iindex(\agent)$ and its component factors.
\end{theorem}

\begin{proof}
Decompose $\Iindex = (\Acon/\Aact) \cdot (\Tcon/\Tcon^{\mathrm{ref}}) \cdot \kappa$ into
three independent factors. Each factor admits depression, normality, or elevation.
Restricting to single-factor deviations (multi-factor failures are superpositions),
there are exactly six single-deviation patterns:

\begin{enumerate}[label=\textbf{(P\arabic*)}]
\item \emph{$\Acon/\Aact$ depressed}: the \textbf{construction-deficient} phenotype.
  Construction phase fails to reach full activity budget. Novel cells are constructed at
  reduced rate or with lower precision.

\item \emph{$\Acon/\Aact$ elevated above reference}: the \textbf{hyper-constructive}
  phenotype. Construction phase exceeds healthy budget; novel cells are generated but
  inadequately tested. Predicted in agents whose perceptual gating fails during
  construction.

\item \emph{$\Tcon/\Tcon^{\mathrm{ref}}$ depressed}: the \textbf{construction-deprived}
  phenotype. Insufficient construction time; too few cells are constructed per unit total
  time even when construction-phase activity is intact.

\item \emph{$\kappa$ depressed}: the \textbf{perceptually decoupled} phenotype. The
  agent constructs cells but cannot test condition~(iv) because perceptual coupling to
  other agents is degraded. Novel cells become private; intelligence collapses to
  idiosyncrasy.

\item \emph{$\Aact$ depressed without proportional $\Acon$ depression}: the
  \textbf{action-cycle disrupted} phenotype. The cognitive cycle is preserved but action
  commitment fails. The agent constructs and tests cells internally but cannot externalise.

\item \emph{$\Acon$ depressed without proportional $\Aact$ depression}: the
  \textbf{construction-cycle disrupted} phenotype. Action commitment is intact but
  framework extension fails. The agent acts competently within existing $\KFW$ but
  cannot grow it.
\end{enumerate}

Each phenotype is decidable in $O(1)$ given measurements of $\Acon, \Aact, \Tcon, \kappa$.
The classification is complete (covers all single-deviation patterns) and exclusive
(each pattern produces a distinct $\Iindex$ signature).
\end{proof}

\begin{remark}[Substrate neutrality of the phenotypes]
The six phenotypes are substrate-neutral. In biological agents, specific cycle failures
correspond to specific physiological conditions that will be examined in Parts~V and~VI
of this monograph. In computational agents, they correspond to architectural or training
pathologies. In organisational agents, they correspond to institutional dysfunctions.
The framework predicts that the same six phenotypes will appear across substrates because
they are forced by the structure of the intelligence cycle, not by the substrate's
particulars.
\end{remark}

\end{multicols}

%% ============================================================
%% Panel 5 — Intelligence Floor and Failure Phenotypes
%% ============================================================
\begin{figure}[!htbp]
    \centering
    \includegraphics[width=\textwidth]{figures/operational-intelligence/panel_5.png}
    \caption{\textbf{Intelligence Floor and the Six Failure-Mode Phenotypes.}
    (\textbf{A})~Bar chart of $\Iindex$ for the healthy reference (seagreen) and
    phenotypes P1--P6. Phenotypes P1 (construction-deficient), P3 (construction-deprived),
    and P4 (perceptually decoupled) show depressed indices; P2 (hyper-constructive) shows
    elevation without proportional gain; P5 and P6 show distinct cycle-disruption patterns
    (Theorem~\ref{thm:ch05-t7}). (\textbf{B})~Intelligence ceiling $1/\beta$ versus
    $|\KFW|$ on a log scale: the ceiling grows but remains finite at every tested framework
    size, confirming Theorem~\ref{thm:ch05-t6}: no maximally intelligent agent exists.
    (\textbf{C})~Factor signatures for the healthy reference and P1--P4: each phenotype
    has a unique three-factor signature, confirming the classification's distinctness.
    (\textbf{D})~Clustering of $30$ agents per phenotype in
    $(\Acon/\Aact,\,\Tcon/\Tcon^{\mathrm{ref}},\,\kappa)$ factor space: the clusters
    occupy distinct, non-overlapping regions, confirming that the index decomposition is
    operationally diagnostic.}
    \label{fig:ch05-panel5-floor}
\end{figure}

%% ============================================================
%% Block 6 — Validation, connections, conclusion
%% ============================================================
\begin{multicols}{2}

\section{Numerical validation}
\label{sec:ch05-validation}

The validation suite comprises thirty experiments in six clusters covering foundations
(C1), cycle dynamics (C2), the intelligence index (C3), collective intelligence (C4),
the floor and bounds (C5), and failure-mode phenotypes (C6). Each experiment compares a
measured quantity to its closed-form theoretical prediction and reports relative error.

\begin{itemize}
\item \textbf{C1} (E1--E5): floor positivity, cell-disjoint admissibility, receiver
  uncertainty product, construction-phase dispersion, action-phase dispersion.
\item \textbf{C2} (E6--E10): phase alternation necessity, construction-phase activity
  ratio, single-phase agent failures, cycle index positivity.
\item \textbf{C3} (E11--E15): index closed-form verification, sensitivity analysis,
  reference-value prediction.
\item \textbf{C4} (E16--E20): aperture-sharing, phase-locking, federation index
  aggregation.
\item \textbf{C5} (E21--E25): intelligence floor positivity, boundedness, monotonicity
  in $\beta$.
\item \textbf{C6} (E26--E30): six phenotype signatures.
\end{itemize}

\begin{table}[H]\centering\small
\begin{tabular}{lcc}
\toprule
\textbf{Cluster} & \textbf{Max rel.\ err.} & \textbf{Verdict}\\
\midrule
C1 foundations       & $\le 10^{-15}$ & \textsc{Pass}\\
C2 cycle dynamics    & $0.00$ (bool.) & \textsc{Pass}\\
C3 index             & $\le 10^{-15}$ & \textsc{Pass}\\
C4 collective        & $0.00$ (bool.) & \textsc{Pass}\\
C5 floor/bounds      & $\le 10^{-15}$ & \textsc{Pass}\\
C6 failure modes     & $0.00$ (bool.) & \textsc{Pass}\\
\midrule
\textbf{Aggregate}   & $\mathbf{\le 10^{-15}}$ & \textbf{30/30}\\
\bottomrule
\end{tabular}
\caption{Validation summary across $30$ experiments. All claims verified to machine
precision.}
\label{tab:ch05-validation}
\end{table}

\section{Connections to existing frameworks}
\label{sec:ch05-connections}

\subsection{Existing intelligence definitions}

The Legg--Hutter definition~\citep{legghutter2007,hutter2005} equates intelligence with
expected reward across a universal class of environments. This is non-operational: it
requires evaluating agents on environments drawn from an uncomputable prior. The index
$\Iindex(\agent)$ is computable directly from agent telemetry.

The Turing test~\citep{turing1950} is substrate-anchored to human behaviour. Our index
is substrate-neutral: the intelligence cycle is defined without reference to any
particular kind of agent.

Bounded-rationality frameworks~\citep{simon1996,russellnorvig2020} share the
bounded-resources axiom but do not address knowledge framework growth, which decision
theory omits by assumption. Growth is the operative component of intelligence that
separates it from competence.

\subsection{Foundational logic}

The receiver floor $\beta > 0$ is the operational analogue of G\"{o}del
incompleteness~\citep{godel1931,tarski1933}, Church's thesis~\citep{church1936}, and
Chaitin complexity~\citep{chaitin1974}: in every case, a strictly positive lower bound on
what can be represented within a finite consistent system. The intelligence cycle adds
the dynamical content absent from these static incompleteness results: the floor is
reduced over time on novel cells, but never eliminated.

\subsection{Uncertainty inequalities and information theory}

The Receiver Uncertainty Principle $\sigma_K \sigma_Y \ge \hbarA$ has the same product
form as the Heisenberg principle~\citep{heisenberg1927,robertson1929} and the
Cram\'er--Rao bound~\citep{rao1945,cramer1946}, with $\hbarA = \beta\tau$ playing the
role of $\hbar$. The index factors carry information-theoretic interpretations:
$\Acon/\Aact$ resembles a Shannon activity ratio~\citep{shannon1948,coverthomas2006};
$\kappa$ resembles a Kullback--Leibler coupling coefficient~\citep{kullback1951};
$\Tcon/\Tcon^{\mathrm{ref}}$ captures temporal allocation in the spirit of Fisher
sufficiency~\citep{fisher1925}.

\subsection{Synchronisation and distributed computing}

The collective intelligence theorem builds on the Kuramoto model and its rigorous
analysis~\citep{kuramoto1984,strogatz2000,acebron2005,pikovsky2001}. The
aperture-sharing condition is the agent-level analogue of phase identity. Common-Cell
Convergence (Chapter~\ref{ch:purpose}) dispenses with iterated common-knowledge
prerequisites~\citep{aumann1976,fagin1995,lewis1969} for coordination.

The Fischer--Lynch--Paterson impossibility~\citep{fischer1985} and the Byzantine
generals problem~\citep{lamport1982} bound consensus in turbulent and cascade regimes.
In the phase-locked regime, partition extinction eliminates the need for explicit
consensus and these impossibility theorems do not apply.

\section{Conclusion}
\label{sec:ch05-conclusion}

Seven theorems produce a substrate-neutral operational definition of intelligence for
any bounded agent. The principal results are:

\begin{enumerate}
\item Intelligence is the capacity for cell-disjoint extension of the knowledge framework
  that reduces effective floor on novel cells and enables synchronisation on those cells.
\item Extension is admissible iff cell-disjoint (T1).
\item Extension occurs only during construction phases by the Receiver Uncertainty
  Principle (T2).
\item Intelligence requires alternation between construction and action phases (T3).
\item The intelligence index $\Iindex(\agent)$ is computable in closed form (T4).
\item Collective intelligence is synchronisation via aperture-sharing isomorphisms (T5).
\item Every realised agent has finite $\Iindex$ --- no maximally intelligent agent exists (T6).
\item Cycle failures admit exactly six distinct phenotypes (T7).
\end{enumerate}

The framework subsumes the results of the preceding chapters --- cell-truth, partition
extinction, catalytic composition, federation entropy, and Lagrangian trajectory
compatibility --- as special cases. It applies uniformly to biological, computational,
and organisational agents.

The six failure-mode phenotypes will recur throughout this monograph: in Part~V as the
molecular Maxwell demons whose activity shifts biological agents between phenotypes, in
Part~VI as computational architectures that can or cannot realise intelligent behaviour,
and in the pharmacological framework of Part~V where specific interventions are
characterised by which phenotype they induce or repair. The substrate-neutral taxonomy
established here provides the diagnostic vocabulary for all of those analyses.

\end{multicols}
